\DTMsavedate{mydate}{2022-06-30}
\maketitle{Computer Engineering \Romannum{1}}{Electronic Engineering}{\DTMusedate{mydate}}

% ----------------------------------------------------- %
% COURSE INFO
\subsection*{Course Information}\label{course:tee2232}
\vspace{0ex}%
\noindent\parbox{0.5\textwidth}{%
    \noindent\begin{tabular}{@{}ll}
                 \textsf{Course:}          & 	Computer Engineering \Romannum{1}    \\
                 \textsf{Course Code:}         & 	TEE 2232    \\
                 \textsf{Credit Hours:}    & 	10
    \end{tabular}}
\vspace{2ex}\hrule\vspace{2ex}

% ----------------------------------------------------- %
% INSTRUCTOR INFO
\subsection*{Lecturer Information}
\noindent\parbox{0.5\textwidth}{%
    \noindent\begin{tabular}{@{}ll}
                 \textsf{Name:}         &     Lovemore Gunda         \\
%\textsf{Phone:}        &    \phone          \\
\textsf{Email:}        &    loveugunda@gmail.com \\
\textsf{Qualifications:}        &  MEng in Electronic Engineering, (Hon) BEng in Electronic Engineering
    \end{tabular}}
\vspace{2ex}\hrule\vspace{2ex}


% ----------------------------------------------------- %
% COURSE DESCRIPTION
\subsection*{Course Synopsis}
Computer components: types of computer systems, system features, function of central processing unit, the motherboard, the expansion bus, system memory map, volatile/non-volatile memory storage methods.
Introduction to operating systems: MSDOS, UNIX, and Windows. Software development, software life cycles, structured programming, introduction to modular design: C-constructs, introduction to programming in C. Packages: computer maintenance and  diagnostic, spreadsheets
\vspace{2ex} \hrule\vspace{2ex}

% ----------------------------------------------------- %
% LEARNING OBJECTIVES
\subsection*{Learning Objectives}
At the end of the course students should be able to:
\begin{outline}
    \1      To identify and describe the functions of computer hardware and computer systems.

    \1  	To understand the basics and administration of various operating systems including DOS, Windows, and Unix/Linux

    \1  	To explain the structure, operation, programming, and applications of computers.

    \1  	To understand maintenance and diagnostics for computers.

\end{outline}
\vspace{2ex}\hrule\vspace{2ex}

% ----------------------------------------------------- %
% COURSE TOPICS
\subsection*{Topics}
\begin{outline}
    \1 Introduction to computer hardware
    \1 Components of a computer
    \1 Introduction to operating systems
    \1 Software development
    \1 Introduction to programming with C
    \1 Computer software packages
    \1 Computer maintenance and diagnostics
\end{outline}
\vspace{2ex} \hrule\vspace{2ex}

% ----------------------------------------------------- %
% Students Assenssment
\subsection*{Assessment of Students}
\begin{outline}
    \1 Students are assessed through:
    \begin{enumerate}
        \item Written assignments.
        \item Practical assignments.
        \item Written in-class test
        \item Practical in-class test
        \item End of semester final test
    \end{enumerate}

\end{outline}
\vspace{2ex}\hrule\vspace{2ex}

% ----------------------------------------------------- %
% Grade Evaluation
\subsection*{Course Evaluation and Grading Policies}
Grades are based on:
Written assignment (\qty{5}{\percent}),
Practical assignment (\qty{5}{\percent}),
Written in-class test (\qty{5}{\percent}),
Practical in-class test (\qty{10}{\percent}),
Final Exam (\qty{75}{\percent})
\vspace{2ex}\hrule\vspace{2ex}

% ----------------------------------------------------- %
% EQUIPMENT
% https://sites.udel.edu/computing-purchases/personal-specs/
% \subsection*{Essential Equipment}

% \vspace{2ex}\hrule\vspace{2ex}

% Policies and Other information
\subsection*{Policies and Other Information}
% Conduct
\subsection*{Key Text}
\begin{itemize}
    \item Operating Systems Design and Implementation, 3rd Edition, 2006 by A. Tanenbaum and A. Woodhull
    \item Modern Operating Systems, 2nd Edition, 2001 by A. Tanenbaum
    \item Software Architecture in Practice, 2nd Edition by L. Bass, P. Clements, R. Kazman
    \item Computing science, 2nd Edition by P. Bishop
\end{itemize}
\vspace{2ex}\hrule
\vspace{2ex}\hrule\vspace{2ex}